\documentclass[11pt]{article}

\usepackage[spanish]{babel}
\usepackage[utf8]{inputenc}
\usepackage[T1]{fontenc}
\usepackage{amsmath, amssymb, amsthm, mathtools}
\usepackage{geometry}
\usepackage{xcolor}
\usepackage{tcolorbox}

\geometry{letterpaper, margin=2.2cm}

% ==============================================================================
% DEFINICIONES DE ENTORNOS PERSONALIZADOS PARA COMPILACIÓN
% ==============================================================================

\newenvironment{motivacion}{%
  \section*{1. Motivación}%
}{}

\newenvironment{preguntaguia}{%
  \begin{tcolorbox}[colback=cyan!5,colframe=cyan!75!black,title=Pregunta Guía]%
}{%
  \end{tcolorbox}%
}

\newenvironment{teoria}{%
  \section*{2. Definición y Teoría}%
}{}

\newenvironment{definicion}[1]{%
  \begin{tcolorbox}[colback=blue!5,colframe=blue!75!black,title=Definición: #1]%
}{%
  \end{tcolorbox}%
}

\newenvironment{teorema}[1]{%
  \begin{tcolorbox}[colback=green!5,colframe=green!75!black,title=Teorema: #1]%
}{%
  \end{tcolorbox}%
}

\newenvironment{alerta}[1]{%
  \begin{tcolorbox}[colback=yellow!10,colframe=orange!80!black,title=Alerta: #1]%
}{%
  \end{tcolorbox}%
}

\newenvironment{procesamiento}[1]{%
  \begin{tcolorbox}[colback=gray!5,colframe=gray!60!black,title=Procedimiento: #1]%
}{%
  \end{tcolorbox}%
}

\newenvironment{aplicacion}{%
  \section*{3. Aplicación y Práctica}%
}{}

\newenvironment{ejemplo}[1]{%
  \begin{tcolorbox}[colback=white,colframe=gray!40,title=Ejemplo: #1]%
}{%
  \end{tcolorbox}%
}

% Comandos para preguntas interactivas
\newenvironment{preguntaalternativas}[1]{\subsection*{Pregunta: #1}\par\vspace{0.1cm}}{\vspace{0.3cm}}
\newcommand{\opcion}[3]{\par\noindent $\bullet$ \textbf{#1} \textit{(#2)} - #3\par\vspace{0.15cm}}

\newenvironment{preguntaverdaderofalso}[1]{\subsection*{Pregunta V/F: #1}\par\vspace{0.1cm}}{\vspace{0.3cm}}
\newcommand{\verdaderofalso}[3]{\par\noindent\textbf{Respuesta correcta: #1} \\ \textit{Si marca Falso:} #2 \\ \textit{Si marca Verdadero:} #3\par\vspace{0.15cm}}

\newenvironment{preguntacasillas}[1]{\subsection*{Selección Múltiple: #1}\par\vspace{0.1cm}}{\vspace{0.3cm}}
\newcommand{\casilla}[3]{\par\noindent $\square$ \textbf{#1} \textit{(#2)} - #3\par\vspace{0.15cm}}

% Entorno Términos Pareados (3 Columnas)
\newenvironment{pareadostrescolumnas}[1]{\subsection*{Términos Pareados: #1}}{\vspace{0.3cm}}
\newcommand{\columnaI}[1]{\textbf{Columna 1:}\begin{enumerate}#1\end{enumerate}}
\newcommand{\columnaII}[1]{\textbf{Columna 2:}\begin{itemize}#1\end{itemize}}
\newcommand{\columnaIII}[1]{\textbf{Columna 3:}\begin{itemize}#1\end{itemize}}
\newcommand{\pareotres}[2]{\textbf{Cruce #1:} #2\par}

% Entornos para ejercicios
\newenvironment{ejercicios}{%
  \section*{4. Ejercicios Resueltos y Propuestos}%
}{}

\newenvironment{ejercicioresuelto}[2]{\subsection*{Ejercicio Resuelto: #1 \small{[#2]}}}{}
\newenvironment{ejerciciodemostracion}[2]{\subsection*{Demostración: #1 \small{[#2]}}}{}
\newenvironment{ejerciciopropuesto}[2]{\subsection*{Ejercicio Propuesto: #1 \small{[#2]}}}{}

\newcommand{\enunciado}[1]{\textbf{Enunciado:} #1\par\vspace{0.2cm}}
\newcommand{\solucion}[1]{\textbf{Solución:} #1\par\vspace{0.4cm}}
\newcommand{\demostracion}[1]{\textbf{Demostración:} #1\par\vspace{0.4cm}}
\newcommand{\pista}[1]{\textbf{Pista:} #1\par\vspace{0.4cm}}

% Entorno Fórmulas
\newenvironment{formulas}{\section*{5. Fórmulas Clave}\begin{description}}{\end{description}}
\newcommand{\formula}[3]{\item[#1:] $#2$ \\ \textit{#3} \vspace{0.2cm}}

% ==============================================================================
% 1. METADATOS TÉCNICOS DE DESPLIEGUE (COMPLETAR OBLIGATORIAMENTE)
% ==============================================================================
% ID_CURSO: calculo-multivariable
% NUMERO_UNIDAD: 1
% TITULO_UNIDAD: Funciones de Varias Variables
% NUMERO_CAPITULO: 1.3
% TITULO_CAPITULO: Límites de Campos Escalares
% ESTADO_BLOQUEADO: false
% ESTADO_COMPLETADO: false
% ==============================================================================

\begin{document}

% ==============================================================================
% PESTAÑA 1: MOTIVACIÓN (contentMotivation)
% ==============================================================================
\begin{motivacion}
  \textbf{¿Cómo te aproximas a un destino cuando no hay un solo camino?}

  En Cálculo de una variable, acercarse a un punto $x = a$ era un asunto sencillo: solo podías caminar desde la izquierda ($x \to a^-$) o desde la derecha ($x \to a^+$). Si ambos caminos laterales te llevaban a la misma altura, el límite existía.

  Sin embargo, en el plano cartesiano $\mathbb{R}^2$, estar en un punto $(x_0, y_0)$ es como estar en el centro de una plaza abierta. Puedes aproximarte a ese destino desde infinitas direcciones distintas: caminando en línea recta con cualquier pendiente, o siguiendo una parábola, una cúbica o cualquier curva imaginable.

  Para que el límite exista en varias variables, el campo escalar debe tender al mismo valor real $L$ **a lo largo de absolutamente todas las trayectorias posibles**. Basta con que un solo camino te entregue una altura distinta para que el límite no exista.

  \begin{preguntaguia}
    Si evaluamos un límite a lo largo de todas las rectas posibles $y = mx$ y en todas obtenemos el mismo valor $L$, ¿podemos asegurar con 100\% de certeza que el límite del campo escalar existe?
  \end{preguntaguia}
\end{motivacion}


% ==============================================================================
% PESTAÑA 2: DEFINICIÓN Y TEORÍA (contentTheory)
% ==============================================================================
\begin{teoria}
  El estudio de límites en campos escalares establece las bases analíticas del cálculo multivariable. Para abordar el límite de una función en el plano, debemos formalizar primero la noción topológica del punto al cual nos aproximamos.

  % --- DEFINICIÓN: PUNTO DE ACUMULACIÓN EN R^2 ---
  \begin{definicion}{Punto de Acumulación en $\mathbb{R}^2$}
    Un punto $(x_0, y_0) \in \mathbb{R}^2$ es un \textbf{punto de acumulación} de un conjunto $D \subseteq \mathbb{R}^2$ si para todo $\delta > 0$, la bola abierta perforada $B_\delta(x_0,y_0) \setminus \{(x_0,y_0)\}$ contiene al menos un punto de $D$.
    
    Formalmente:
    $$ \forall \delta > 0, \quad \Big( \{ (x,y) \in \mathbb{R}^2 \colon 0 < \sqrt{(x-x_0)^2 + (y-y_0)^2} < \delta \} \Big) \cap D \neq \emptyset $$
    
    \textbf{Recordatorio:} La expresión $\sqrt{(x-x_0)^2 + (y-y_0)^2}$ representa la distancia euclidiana entre los puntos $(x,y)$ y $(x_0,y_0)$, la cual también se puede escribir como la norma de la diferencia $\|(x,y) - (x_0,y_0)\|$.

    \textbf{Ejemplos en $\mathbb{R}^2$:}
    \begin{itemize}
      \item Sea $D = \{ (x,y) \in \mathbb{R}^2 \colon x^2 + y^2 < 1 \}$ (disco abierto de radio 1). Todos los puntos de la frontera $x^2 + y^2 = 1$ son puntos de acumulación de $D$, a pesar de no pertenecer al conjunto $D$.
      \item Sea $D = \mathbb{R}^2 \setminus \{(0,0)\}$. El origen $(0,0)$ no pertenece a $D$, pero es un punto de acumulación de $D$.
    \end{itemize}
  \end{definicion}

  % --- DEFINICIÓN: LÍMITE FORMAL EPSILON-DELTA ---
  \begin{definicion}{Límite de un Campo Escalar ($\varepsilon$-$\delta$)}
    Sea $f \colon D \subseteq \mathbb{R}^2 \to \mathbb{R}$ un campo escalar. Exigimos que $(a, b)$ sea punto de acumulación de $\operatorname{dom}(f)$ porque así garantizamos que el conjunto de puntos $(x,y)$ tales que la distancia entre $(x,y)$ y $(a,b)$ es menor a $\delta$, no es vacío, y por tanto la aproximación tiene sentido.
    
    Decimos que el \textbf{límite} de $f(x,y)$ cuando $(x,y)$ tiende a $(a, b)$ es $\ell \in \mathbb{R}$, denotado por:
    $$ \lim_{(x,y) \to (a, b)} f(x,y) = \ell $$
    si y solo si:
    $$ \forall \varepsilon > 0, \; \exists \delta > 0 \text{ tal que } 0 < \overbrace{\underbrace{\|(x,y)-(a,b)\|}_{\begin{array}{l} \text{distancia entre} \\ (x,y) \text{ y } (a,b) \end{array}}}^{=\sqrt{(x-a)^2 + (y-b)^2}} < \delta \Longrightarrow \underbrace{|f(x,y)-\ell|}_{\begin{array}{l} \text{distancia entre } \\ f(x,y) \text{ y } \ell \end{array}} < \varepsilon $$
  \end{definicion}

  % --- TEOREMA: UNICIDAD DEL LÍMITE ---
  \begin{teorema}{Unicidad del Límite}
    Sea $f \colon D \subseteq \mathbb{R}^2 \to \mathbb{R}$ un campo escalar y sea $(x_0, y_0)$ un punto de acumulación de $\operatorname{dom}(f)$. Si el límite de $f(x,y)$ cuando $(x,y) \to (x_0, y_0)$ existe, entonces este valor es único. Es decir, si:
    $$ \lim_{(x,y) \to (x_0, y_0)} f(x,y) = L_1 \quad \wedge \quad \lim_{(x,y) \to (x_0, y_0)} f(x,y) = L_2 $$
    entonces $L_1 = L_2$.
  \end{teorema}

  \begin{proof}[\textbf{Demostración:}]
    Procederemos por contradicción. Supongamos que $L_1 \neq L_2$, de modo que $|L_1 - L_2| > 0$. Definamos $\varepsilon = \dfrac{|L_1 - L_2|}{2} > 0$.

    Por la definición formal de límite:
    \begin{itemize}
      \item Para $\displaystyle \lim_{(x,y) \to (x_0,y_0)} f(x,y) = L_1$, existe $\delta_1 > 0$ tal que si $0 < \sqrt{(x-x_0)^2 + (y-y_0)^2} < \delta_1$, entonces $|f(x,y) - L_1| < \varepsilon$.
      \item Para $\displaystyle \lim_{(x,y) \to (x_0,y_0)} f(x,y) = L_2$, existe $\delta_2 > 0$ tal que si $0 < \sqrt{(x-x_0)^2 + (y-y_0)^2} < \delta_2$, entonces $|f(x,y) - L_2| < \varepsilon$.
    \end{itemize}

    Dado que $(x_0,y_0)$ es punto de acumulación de $\operatorname{dom}(f)$, en la bola perforada de radio $\delta = \min(\delta_1, \delta_2)$ existe al menos un punto $(x,y) \in \operatorname{dom}(f)$. Para dicho punto, aplicando la desigualdad triangular:
    $$ |L_1 - L_2| = |(L_1 - f(x,y)) + (f(x,y) - L_2)| \leq |f(x,y) - L_1| + |f(x,y) - L_2| $$
    Sustituyendo los acotamientos:
    $$ |L_1 - L_2| < \varepsilon + \varepsilon = 2\varepsilon = 2 \left( \dfrac{|L_1 - L_2|}{2} \right) = |L_1 - L_2| $$
    Llegamos a la contradicción estricta $|L_1 - L_2| < |L_1 - L_2|$. Por lo tanto, el supuesto inicial es falso y se concluye que $L_1 = L_2$.
  \end{proof}

  % --- EJEMPLO DEMOSTRACIÓN POR DEFINICIÓN ---
  \begin{ejemplo}{Cálculo de Límite por Definición $\varepsilon$-$\delta$}
    Demuestre por definición que:
    $$ \lim_{(x,y) \to (0,0)} \dfrac{x^2}{\sqrt{x^2 + y^2}} = 0 $$
    
    \textbf{Desarrollo:}
    Dado $\varepsilon > 0$, debemos encontrar $\delta > 0$ tal que si $0 < \sqrt{x^2 + y^2} < \delta$, entonces:
    $$ \left| \dfrac{x^2}{\sqrt{x^2 + y^2}} - 0 \right| < \varepsilon $$
    
    Como $x^2 \leq x^2 + y^2$ para todo $(x,y) \in \mathbb{R}^2$, acotamos la expresión:
    $$ \left| \dfrac{x^2}{\sqrt{x^2 + y^2}} \right| = \dfrac{x^2}{\sqrt{x^2 + y^2}} \leq \dfrac{x^2 + y^2}{\sqrt{x^2 + y^2}} = \sqrt{x^2 + y^2} $$
    
    Por lo tanto, basta elegir $\delta = \varepsilon$. De esta forma, si $0 < \sqrt{x^2 + y^2} < \delta$, se cumple directamente que:
    $$ \left| \dfrac{x^2}{\sqrt{x^2 + y^2}} - 0 \right| \leq \sqrt{x^2 + y^2} < \delta = \varepsilon $$
    Quedando demostrado el límite por definición.
  \end{ejemplo}

  En general, es complicado estudiar un límite usando la definición $\varepsilon$-$\delta$. Las siguientes herramientas analíticas nos permiten calcular límites de forma directa:

  % --- TEOREMA: ÁLGEBRA DE LÍMITES ---
  \begin{teorema}{Álgebra de Límites}
    Si $\displaystyle \lim_{(x,y) \to (a,b)} f(x,y) = L$ y $\displaystyle \lim_{(x,y) \to (a,b)} g(x,y) = M$, y $c \in \mathbb{R}$ es una constante, entonces se cumplen las siguientes propiedades algebraicas:
    \begin{enumerate}
      \item \textbf{Múltiplo escalar:} $\displaystyle \lim_{(x,y) \to (a,b)} [c \cdot f(x,y)] = c \cdot L$
      \item \textbf{Suma y Resta:} $\displaystyle \lim_{(x,y) \to (a,b)} [f(x,y) \pm g(x,y)] = L \pm M$
      \item \textbf{Producto:} $\displaystyle \lim_{(x,y) \to (a,b)} [f(x,y) \cdot g(x,y)] = L \cdot M$
      \item \textbf{Cociente:} $\displaystyle \lim_{(x,y) \to (a,b)} \left[ \dfrac{f(x,y)}{g(x,y)} \right] = \dfrac{L}{M}$, siempre que $M \neq 0$.
    \end{enumerate}
  \end{teorema}

  % --- RESULTADO 1: LÍMITES NOTABLES ---
  \begin{procesamiento}{Límites Notables}
    Los límites notables de Cálculo de una Variable sirven para calcular límites de funciones de Varias Variables mediante sustituciones algebraicas directas.
    
    \textbf{Ejemplo:}
    $$ \lim_{(x,y) \to (0,0)} \dfrac{\sin(xy)}{xy} $$
    Haciendo el cambio de variable $u = xy$, observamos que cuando $(x,y) \to (0,0)$, la variable escalar $u \to 0$. Aplicando el límite notable conocido en una variable:
    $$ \lim_{(x,y) \to (0,0)} \dfrac{\sin(xy)}{xy} = \lim_{u \to 0} \dfrac{\sin(u)}{u} = 1 $$
  \end{procesamiento}

  % --- RESULTADO 2: TEOREMA DE CERO POR ACOTADO ---
  \begin{teorema}{Teorema de Cero por Acotado}
    Si $\displaystyle \lim_{(x,y) \to (a,b)} f(x,y) = 0$ y $g(x,y)$ es una función acotada en una vecindad perforada alrededor de $(a,b)$, entonces:
    $$ \lim_{(x,y) \to (a,b)} [f(x,y)g(x,y)] = 0 $$
  \end{teorema}

  \begin{proof}[\textbf{Demostración:}]
    Como $g(x,y)$ es acotada, existe una constante $M > 0$ tal que $|g(x,y)| \leq M$ en una bola perforada de radio $\delta_1 > 0$ centrada en $(a,b)$.
    
    Por otro lado, como $\displaystyle \lim_{(x,y) \to (a,b)} f(x,y) = 0$, para todo $\varepsilon > 0$ existe $\delta_2 > 0$ tal que si $0 < \sqrt{(x-a)^2 + (y-b)^2} < \delta_2$, entonces $|f(x,y)| < \dfrac{\varepsilon}{M}$.
    
    Eligiendo $\delta = \min(\delta_1, \delta_2)$, si $0 < \sqrt{(x-a)^2 + (y-b)^2} < \delta$, tenemos:
    $$ |f(x,y)g(x,y) - 0| = |f(x,y)| \cdot |g(x,y)| < \dfrac{\varepsilon}{M} \cdot M = \varepsilon $$
    Por lo tanto, $\displaystyle \lim_{(x,y) \to (a,b)} [f(x,y)g(x,y)] = 0$.
  \end{proof}

  % --- ALERTA: PROPIEDADES QUE DEJAN DE SER VÁLIDAS ---
  \begin{alerta}{Propiedades que dejan de ser válidas}
    En Cálculo de varias variables \textbf{no son aplicables}:
    \begin{itemize}
      \item \textbf{Límites laterales simples:} No bastan dos direcciones (izquierda/derecha). En $\mathbb{R}^2$ la aproximación a un punto $(a,b)$ se puede realizar desde infinitas trayectorias planas.
      \item \textbf{Regla de L'Hôpital:} No existe una extensión directa a funciones de varias variables, debido a la ausencia de un cociente de derivadas unidimensionales con sentido geométrico único.
    \end{itemize}
  \end{alerta}

\end{teoria}


% ==============================================================================
% PESTAÑA 3: APLICACIÓN Y PRÁCTICA (contentApplication)
% ==============================================================================
\begin{aplicacion}
  A continuación, aplicaremos la definición formal $\varepsilon$-$\delta$, los límites notables y los criterios de acotamiento desarrollados en la teoría para calcular rigurosamente límites multivariables.

  % --- EJEMPLO DESARROLLADO ---
  \begin{ejemplo}{Demostración por Definición $\varepsilon$-$\delta$}
    Calcule por definición el siguiente límite:
    $$ \lim_{(x,y) \to (0,0)} \dfrac{x^2 y}{x^2 + y^2} $$

    \textbf{Solución y Desarrollo Riguroso:}
    
    \textbf{1. Candidato a Límite:}
    Evaluando formalmente a lo largo de la recta $y = x$:
    $$ \lim_{x \to 0} \dfrac{x^2 (x)}{x^2 + x^2} = \lim_{x \to 0} \dfrac{x^3}{2x^2} = \lim_{x \to 0} \dfrac{x}{2} = 0 $$
    Sospechamos fuertemente que el límite existe y su valor es $L = 0$.

    \textbf{2. Demostración mediante la definición $\varepsilon$-$\delta$:}
    Debemos probar que para todo $\varepsilon > 0$, existe $\delta > 0$ tal que si $0 < \sqrt{x^2 + y^2} < \delta$, entonces:
    $$ \left| \dfrac{x^2 y}{x^2 + y^2} - 0 \right| < \varepsilon $$

    \textbf{3. Acotamiento algebraico:}
    Observamos que para todo $(x,y) \neq (0,0)$, se cumple la desigualdad $x^2 \leq x^2 + y^2$, de donde se deduce que:
    $$ \dfrac{x^2}{x^2 + y^2} \leq 1 $$

    Multiplicando por $|y|$ a ambos lados:
    $$ \left| \dfrac{x^2 y}{x^2 + y^2} \right| = \dfrac{x^2}{x^2 + y^2} \cdot |y| \leq 1 \cdot |y| = |y| $$

    Además, sabemos que $|y| = \sqrt{y^2} \leq \sqrt{x^2 + y^2}$. Por lo tanto:
    $$ \left| \dfrac{x^2 y}{x^2 + y^2} - 0 \right| \leq \sqrt{x^2 + y^2} $$

    \textbf{4. Elección del $\delta$:}
    Dado $\varepsilon > 0$, elegimos $\delta = \varepsilon$. Si $0 < \sqrt{x^2 + y^2} < \delta$, tenemos:
    $$ \left| \dfrac{x^2 y}{x^2 + y^2} - 0 \right| \leq \sqrt{x^2 + y^2} < \delta = \varepsilon $$

    Queda demostrado formalmente que:
    $$ \lim_{(x,y) \to (0,0)} \dfrac{x^2 y}{x^2 + y^2} = 0 $$
  \end{ejemplo}

  % --- TIPO 1: PREGUNTA DE ALTERNATIVAS ---
  \begin{preguntaalternativas}{Aplicación de Teoremas de Acotamiento}
    Para calcular $\displaystyle \lim_{(x,y) \to (0,0)} y^2 \cos\left(\dfrac{1}{xy}\right)$, ¿cuál es el argumento analítico correcto?
    
    \opcion{El límite no existe porque la función coseno oscila infinitamente cuando su argumento crece sin cota.}{Incorrecto}{Aunque es cierto que el coseno oscila, estás olvidando analizar qué ocurre con el factor $y^2$ que lo multiplica.}
    \opcion{Se debe utilizar la Regla de L'Hôpital derivando parcialmente respecto a $y$.}{Incorrecto}{Recuerda la advertencia teórica: la Regla de L'Hôpital no es aplicable directamente en funciones de varias variables.}
    \opcion{El límite es $0$ aplicando el Teorema de Cero por Acotado, ya que $\displaystyle \lim_{(x,y)\to(0,0)} y^2 = 0$ y el coseno siempre está acotado entre $-1$ y $1$.}{Correcto}{¡Excelente! Tienes un término que tiende a cero ($y^2$) multiplicado por una función acotada ($|\cos(u)| \leq 1$). Esto garantiza que el producto tiende a cero.}
  \end{preguntaalternativas}

  % --- TIPO 2: PREGUNTA DE VERDADERO O FALSO ---
  \begin{preguntaverdaderofalso}{Unicidad del Límite y Puntos de Acumulación}
    Si evaluamos un campo escalar $f(x,y)$ acercándonos al origen a lo largo de los ejes $x=0$ e $y=0$, y en ambos casos el límite es $L=5$, entonces el Teorema de Unicidad garantiza que $\displaystyle \lim_{(x,y) \to (0,0)} f(x,y) = 5$.
    
    \verdaderofalso{F}
      {¡Correcto! El Teorema de Unicidad afirma que \textbf{si el límite general existe}, es único. Aproximarse por solo dos caminos no garantiza la existencia del límite general.}
      {Incorrecto. Cuidado con esta trampa lógica. Evaluar dos caminos y obtener el mismo resultado no demuestra la existencia del límite general. La unicidad se aplica solo \textbf{después} de haber probado que el límite existe para toda trayectoria y vecindad.}
  \end{preguntaverdaderofalso}

  % --- TIPO 3: PREGUNTA DE CASILLAS ---
  \begin{preguntacasillas}{Elementos de la Definición Formal}
    Respecto a la definición formal de límite $\varepsilon$-$\delta$ en $\mathbb{R}^2$, selecciona \textbf{todas} las afirmaciones que sean matemáticamente correctas:
    
    \casilla{El punto $(a, b)$ debe pertenecer obligatoriamente a $\operatorname{dom}(f)$.}{Incorrecto}{Falso. Solo se exige que $(a,b)$ sea un punto de acumulación de $\operatorname{dom}(f)$, no que pertenezca a él.}
    \casilla{El valor de $\delta$ generalmente depende de la tolerancia $\varepsilon$ elegida o entregada.}{Correcto}{¡Correcto! En la demostración, construimos $\delta$ en función de $\varepsilon$ (por ejemplo, $\delta = \varepsilon$) para forzar el acotamiento.}
    \casilla{La inecuación $\sqrt{(x-a)^2 + (y-b)^2} < \delta$ describe geométricamente una bola abierta (disco) en el plano.}{Correcto}{¡Exacto! Esa es la definición de distancia euclidiana delimitando una región circular en $\mathbb{R}^2$.}
    \casilla{El Teorema del Sándwich (Acotamiento) es una consecuencia directa de esta definición.}{Correcto}{¡Muy bien! Acotar la función entre $-g(x,y)$ y $g(x,y)$ es justamente lo que hacemos en el álgebra de desigualdades para forzar que el resultado sea menor a $\varepsilon$.}
  \end{preguntacasillas}

  % --- TIPO 4: TÉRMINOS PAREADOS (3 Columnas) ---
  \begin{pareadostrescolumnas}{Asociación: Límite, Técnica Analítica y Resultado}
    
    % COLUMNA 1: Límites
    \columnaI{
      \item $\displaystyle \lim_{(x,y)\to(0,0)} \dfrac{\sin(x^2+y^2)}{x^2+y^2}$
      \item $\displaystyle \lim_{(x,y)\to(1,2)} (3xy - y^2)$
      \item $\displaystyle \lim_{(x,y)\to(0,0)} x^2 \arctan\left(\dfrac{1}{y}\right)$
      \item $\displaystyle \lim_{(x,y)\to(3,3)} \dfrac{x^2-y^2}{x-y}$
      \item $\displaystyle \lim_{(x,y)\to(0,0)} \dfrac{1-\cos(x^2+y^2)}{(x^2+y^2)^2}$
    }
    
    % COLUMNA 2: Técnicas - Mezcladas
    \columnaII{
      \item Simplificación algebraica (Factorización)
      \item Teorema de Cero por Acotado
      \item Límite Notable con coseno
      \item Evaluación directa (Álgebra de límites)
      \item Límite Notable trigonométrico (Seno)
    }
    
    % COLUMNA 3: Resultados - Mezclados
    \columnaIII{
      \item $2$
      \item $6$
      \item $0$
      \item $1$
      \item $1/2$
    }

    % --- SOLUCIONES Y RETROALIMENTACIÓN ---
    \pareotres{1-E-IV}{¡Correcto! Usando la sustitución $u = x^2+y^2$, el límite se transforma en el notable $\displaystyle \lim_{u\to0} \dfrac{\sin(u)}{u} = 1$.}
    \pareotres{2-D-I}{¡Excelente! Al ser un polinomio y $(1,2)$ pertenecer al dominio, evaluamos directamente: $3(1)(2) - 2^2 = 6 - 4 = 2$.}
    \pareotres{3-B-III}{¡Perfecto! El término $x^2 \to 0$ y la función arco tangente está estrictamente acotada entre $-\pi/2$ y $\pi/2$. Por el Teorema de Cero por Acotado, el límite es $0$.}
    \pareotres{4-A-II}{¡Muy bien! Se factoriza la diferencia de cuadrados en el numerador como $(x-y)(x+y)$. Al simplificar $(x-y)$, nos queda $\displaystyle \lim (x+y) = 3+3 = 6$.}
    \pareotres{5-C-V}{¡Magistral! Haciendo el cambio $u = x^2+y^2$, obtenemos el límite notable canónico $\displaystyle \lim_{u\to0} \dfrac{1-\cos(u)}{u^2} = \dfrac{1}{2}$.}

  \end{pareadostrescolumnas}

\end{aplicacion}


% ==============================================================================
% PESTAÑA 4: EJERCICIOS RESUELTOS Y PROPUESTOS (contentExercises)
% ==============================================================================
\begin{ejercicios}

  % --- 1. EJERCICIO RESUELTO: Álgebra de Límites y Resultados Previos ---
  \begin{ejercicioresuelto}{Límite vía Separación y Resultados Conocidos}{nivel-2}
    \enunciado{
      A partir de los resultados analíticos previamente calculados, determine el valor de:
      $$ \lim_{(x,y) \to (0,0)} \dfrac{x^2 y + x y^2}{x^2 + y^2} $$
    }
    \solucion{
      \textbf{1. Descomposición por Aditividad:}
      Aprovechando la estructura del numerador, separamos la fracción en la suma de dos términos racionales:
      $$ \dfrac{x^2 y + x y^2}{x^2 + y^2} = \dfrac{x^2 y}{x^2 + y^2} + \dfrac{x y^2}{x^2 + y^2} $$

      \textbf{2. Aplicación de Resultados Previos:}
      En la Pestaña 3 de Aplicación demostramos por definición $\varepsilon$-$\delta$ que:
      $$ \lim_{(x,y) \to (0,0)} \dfrac{x^2 y}{x^2 + y^2} = 0 $$

      Por simetría algebraica de las variables $x$ e $y$ (o bien aplicando un acotamiento idéntico $\left|\dfrac{x y^2}{x^2 + y^2}\right| \leq |x| \leq \sqrt{x^2+y^2}$), se verifica directamente que:
      $$ \lim_{(x,y) \to (0,0)} \dfrac{x y^2}{x^2 + y^2} = 0 $$

      \textbf{3. Conclusión por el Álgebra de Límites:}
      Dado que ambos límites individuales existen y son finitos, aplicamos la propiedad de la suma del Álgebra de Límites:
      \[\begin{aligned}
        \lim_{(x,y) \to (0,0)} \dfrac{x^2 y + x y^2}{x^2 + y^2} &= \lim_{(x,y) \to (0,0)} \dfrac{x^2 y}{x^2 + y^2} + \lim_{(x,y) \to (0,0)} \dfrac{x y^2}{x^2 + y^2} \\[0.5em]
        &= 0 + 0 = 0
      \end{aligned}\]
    }
  \end{ejercicioresuelto}

  % --- 2. EJERCICIO DE DEMOSTRACIÓN ---
  \begin{ejerciciodemostracion}{Invarianza por Unicidad y Construcción $\varepsilon$-$\delta$}{nivel-3}
    \enunciado{
      Demuestre formalmente que si $\displaystyle \lim_{(x,y) \to (a,b)} f(x,y) = L$, entonces $\displaystyle \lim_{(x,y) \to (a,b)} |f(x,y)| = |L|$.
    }
    \demostracion{
      Sea $\varepsilon > 0$ arbitrario. Dado que $\displaystyle \lim_{(x,y) \to (a,b)} f(x,y) = L$, existe un $\delta > 0$ tal que si $0 < \sqrt{(x-a)^2 + (y-b)^2} < \delta$, entonces:
      $$ |f(x,y) - L| < \varepsilon $$

      Por la propiedad de la desigualdad triangular en su forma reversa para números reales, sabemos que para cualesquiera $A, B \in \mathbb{R}$:
      $$ ||A| - |B|| \leq |A - B| $$

      Haciendo $A = f(x,y)$ y $B = L$, tenemos:
      $$ ||f(x,y)| - |L|| \leq |f(x,y) - L| $$

      Por lo tanto, para el mismo $\delta > 0$, si $0 < \sqrt{(x-a)^2 + (y-b)^2} < \delta$, se cumple encadenadamente que:
      $$ ||f(x,y)| - |L|| \leq |f(x,y) - L| < \varepsilon $$

      Quedando demostrado por definición que $\displaystyle \lim_{(x,y) \to (a,b)} |f(x,y)| = |L|$.
    }
  \end{ejerciciodemostracion}

  % --- 3. EJERCICIO PROPUESTO 1: Combinación de Seno y Raíz ---
  \begin{ejerciciopropuesto}{Límite Combinado de Seno y Raíz}{nivel-2}
    \enunciado{
      Evalúe analíticamente el siguiente límite demostrando su convergencia:
      $$ \lim_{(x,y) \to (0,0)} \dfrac{x^2 \sin(x^2 + y^2)}{\sqrt{x^2 + y^2}} $$
    }
    \pista{
      Puedes reescribir la expresión reordenando sus términos como $x^2 \cdot \dfrac{\sin(x^2 + y^2)}{x^2 + y^2} \cdot \sqrt{x^2 + y^2}$. Utiliza el límite notable $\displaystyle \lim_{u \to 0} \dfrac{\sin(u)}{u} = 1$ haciendo $u = x^2 + y^2$ junto con el Álgebra de Límites, o bien acota utilizando que $|\sin(u)| \leq 1$ para aplicar Cero por Acotado.
    }
  \end{ejerciciopropuesto}

  % --- 4. EJERCICIO PROPUESTO 2: Factorización ---
  \begin{ejerciciopropuesto}{Resolución de Indeterminaciones por Factorización}{nivel-1}
    \enunciado{
      Evalúe el siguiente límite analítico determinando la factorización adecuada:
      $$ \lim_{(x,y) \to (3,3)} \dfrac{x^3 - y^3}{x - y} $$
    }
    \pista{
      Aplica la diferencia de cubos $A^3 - B^3 = (A - B)(A^2 + AB + B^2)$ en el numerador y simplifica el factor $(x-y)$.
    }
  \end{ejerciciopropuesto}

  % --- 5. EJERCICIO PROPUESTO 3: Cero por Acotado ---
  \begin{ejerciciopropuesto}{Funciones Oscilantes Acotadas}{nivel-2}
    \enunciado{
      Demuestre la existencia y calcule el valor del límite:
      $$ \lim_{(x,y) \to (0,0)} x^2 \sin\left(\dfrac{1}{xy}\right) $$
    }
    \pista{
      Aplica el Teorema de Cero por Acotado notando que $x^2 \to 0$ y que $\left|\sin\left(\dfrac{1}{xy}\right)\right| \leq 1$.
    }
  \end{ejerciciopropuesto}

\end{ejercicios}


% ==============================================================================
% PESTAÑA 5: FÓRMULAS CLAVE (contentFormulas)
% ==============================================================================
\begin{formulas}

  \formula{Definición Formal de Límite ($\varepsilon$-$\delta$)}
    {\forall \varepsilon > 0, \; \exists \delta > 0 \colon 0 < \sqrt{(x-a)^2 + (y-b)^2} < \delta \implies |f(x,y) - \ell| < \varepsilon}
    {Condición topológica estricta que garantiza que, para cualquier margen de error $\varepsilon$ en la salida $Z$, existe un disco de radio $\delta$ en el plano $XY$ donde todas las imágenes están atrapadas.}

  \formula{Distancia Euclidiana en $\mathbb{R}^2$}
    {d((x,y), (a,b)) = \sqrt{(x-a)^2 + (y-b)^2} = \|(x,y) - (a,b)\|}
    {Medida de la separación entre dos puntos en el plano. Es el núcleo de la inecuación que define el radio $\delta$ en la aproximación al punto de acumulación.}

  \formula{Álgebra de Límites}
    {\begin{aligned}
      \lim [f(x,y) \pm g(x,y)] &= L \pm M \\[0.5em]
      \lim [f(x,y) \cdot g(x,y)] &= L \cdot M \\[0.5em]
      \lim \left[ \dfrac{f(x,y)}{g(x,y)} \right] &= \dfrac{L}{M} \quad (M \neq 0)
    \end{aligned}}
    {Propiedades operativas que permiten separar y evaluar límites complejos por partes, siempre y cuando los límites individuales $L$ y $M$ existan y sean finitos.}

  \formula{Teorema de Cero por Acotado}
    {\lim_{(x,y) \to (a,b)} f(x,y) = 0 \quad \wedge \quad |g(x,y)| \leq M \implies \lim_{(x,y) \to (a,b)} [f(x,y)g(x,y)] = 0}
    {Herramienta analítica para resolver indeterminaciones. Si un término se anula y el otro permanece acotado, el producto tiende indefectiblemente a cero.}

  \formula{Límite Notable Clásico}
    {\lim_{u \to 0} \dfrac{\sin(u)}{u} = 1 \implies \lim_{(x,y) \to (0,0)} \dfrac{\sin(xy)}{xy} = 1}
    {Sustitución directa ($u = xy$ o $u = x^2+y^2$) que permite importar resultados conocidos de Cálculo de una variable al plano bidimensional.}

\end{formulas}

\end{document}